From a step-response test with an input amplitude of $50^\circ\mathrm{C}$
applied to heater 1, the plant dynamics were identified. The feedback signal is
based on sensor 1, while the tracking objective refers to the relative
temperature measured by sensor 2. The Ziegler-Nichols step-response method was
adopted using a \gls{fopdt} model, yielding an approximation for each sensor, as
given by \Eqref{planta-sensor-1-zn,planta-sensor-2-zn}:

\begin{align}
\eqlabel{planta-sensor-1-zn}
G_{1}(s) &= \frac{0.92}{188.03s + 1} \cdot e^{-6.37s}\\
\eqlabel{planta-sensor-2-zn}
G_{2}(s) &= \frac{0.91}{183.94s + 1} \cdot e^{-11.66s}
\end{align}

The sensor 1 model, which was used for controller design, was discretized using
\gls{zoh} with $T_s = \SI{15}{\second}$, resulting in \Eqref{planta-gz-zoh}.

The model was evaluated by comparing the step response of the \gls{fopdt} model with the experimental data, showing good agreement in both the static gain and the identified time constant.